% ============================================================================
%  Versión corta del asistente NL2C para la memoria (2 a 3 páginas).
%  El desarrollo completo está en nl2c.tex; aquí se condensan las cinco partes del
%  formato de la asignatura de algoritmos y las tablas clave.
%  Compilar:  pdflatex nl2c_memoria && bibtex nl2c_memoria && pdflatex nl2c_memoria (x2)
% ============================================================================
\documentclass[11pt,letterpaper]{article}
% Preámbulo del informe NL2C. Sigue las convenciones de docs/algoritmos/preambulo.tex
% (idioma, seudocódigo con palabras clave en español, \LineComment con ▷, teoremas) sin los
% paquetes que no están instalados en esta máquina (csquotes, siunitx, threeparttable, xurl).
\usepackage[T1]{fontenc}
\usepackage[utf8]{inputenc}
\usepackage[spanish,es-noshorthands,es-tabla,es-nodecimaldot]{babel}
\usepackage{lmodern}
\usepackage{microtype}

\usepackage[margin=2.6cm,headheight=14pt]{geometry}
\usepackage{fancyhdr}
\pagestyle{fancy}
\fancyhf{}
\fancyhead[L]{\small\nouppercase{\leftmark}}
\fancyfoot[C]{\small\thepage}
\renewcommand{\headrulewidth}{0.4pt}
\setlength{\parskip}{0.35em}

\usepackage{amsmath,amssymb,amsthm,mathtools}
\usepackage{booktabs,multirow,array}
\newcolumntype{L}[1]{>{\raggedright\arraybackslash}p{#1}}
\usepackage{graphicx}
\usepackage{float}
\usepackage[font=small,labelfont=bf]{caption}
\usepackage[inline]{enumitem}
\setlist{itemsep=0.15em,topsep=0.3em}

\usepackage{algorithm}
\usepackage[noend]{algpseudocode}
\floatname{algorithm}{Algoritmo}
\algrenewcommand\algorithmicrequire{\textbf{Entrada:}}
\algrenewcommand\algorithmicensure{\textbf{Salida:}}
\algrenewcommand\algorithmicif{\textbf{si}}
\algrenewcommand\algorithmicthen{\textbf{entonces}}
\algrenewcommand\algorithmicelse{\textbf{si no}}
\algrenewcommand\algorithmicfor{\textbf{para}}
\algrenewcommand\algorithmicforall{\textbf{para todo}}
\algrenewcommand\algorithmicdo{\textbf{hacer}}
\algrenewcommand\algorithmicwhile{\textbf{mientras}}
\algrenewcommand\algorithmicreturn{\textbf{devolver}}
\algnewcommand{\LineComment}[1]{\State \(\triangleright\) \textit{#1}}
\algrenewcommand\alglinenumber[1]{\scriptsize #1:}

\usepackage{xcolor}
\usepackage[hidelinks,pdfusetitle,bookmarksnumbered]{hyperref}
\usepackage[spanish,nameinlink,noabbrev]{cleveref}

\theoremstyle{definition}
\newtheorem{definicion}{Definición}[section]

% Comillas angulares de la RAE en lugar de csquotes.
\newcommand{\enquote}[1]{«#1»}

\crefname{table}{tabla}{tablas}
\Crefname{table}{Tabla}{Tablas}
\crefname{algorithm}{algoritmo}{algoritmos}
\Crefname{algorithm}{Algoritmo}{Algoritmos}

\usepackage[authoryear,round]{natbib}
\geometry{margin=2.1cm}
\fancyhead[R]{\small Asistente NL2C}
\newcommand{\R}{\mathcal{R}}
\newcommand{\cod}[1]{\texttt{#1}}
\setlength{\parskip}{0.3em}

\begin{document}

\section{Asistente de lenguaje natural a restricciones (NL2C)}

En esta secci\'on se resume el asistente que convierte la frase de un acudiente en la restricci\'on que usa el recomendador. Para ello,
primero se analiza el problema y se definen sus entradas y salidas, despu\'es se describe la estrategia y su seudoc\'odigo, luego se
estima el costo y, al final, se presentan los resultados de la evaluaci\'on y sus l\'imites. El desarrollo completo est\'a en el informe
\cod{nl2c.tex}.

\subsection{An\'alisis del problema}

La propuesta aprobada prometi\'o que el acudiente pudiera describir en espa\~nol lo que quiere, por ejemplo \enquote{Minnie para mi
hija, m\'aximo 500 mil}, en lugar de llenar un formulario de opciones cerradas. A esa tarea se le llama NL2C (del ingl\'es
\emph{natural language to constraints}): transformar un texto en el objeto de restricciones que un programa ya consume. Aqu\'i el
programa es el optimizador NSGA-II \citep{deb2002nsga2}, que necesita filtros, pesos y preferencias, no una frase. Las observaciones
que guían el dise\~no son estas.
\begin{enumerate}[label=\textbf{O\arabic*.},leftmargin=*]
  \item El texto es una cadena $t\in\Sigma^{*}$, donde $\Sigma$ es el alfabeto Unicode; es corta, ambigua y con errores de ortograf\'ia.
  \item La salida es $\R=\langle F,W,P\rangle$ (\cod{filtros\_duros}, \cod{pesos\_NSGA}, \cod{preferencias\_PRM}) con los dominios de
  la \cref{tab:dominios}. Un campo que el acudiente no menciona queda en su valor por omisi\'on, y uno que menciona fuera de
  dominio se lleva al dominio y se le avisa.
  \item $t$ y el texto de la lista escolar los puede escribir un atacante, de modo que ninguno puede cambiar qu\'e acciones ejecuta el sistema
  \citep{greshake2023indirect}. Adem\'as, un error cuesta dinero al acudiente, por lo que nada se aplica sin su confirmaci\'on.
\end{enumerate}
No se resuelve con reglas fijas por tres razones: la frase es ambigua (\enquote{calidad primero y precio despu\'es} fija pesos pero
no el estilo de gasto), los nombres de personajes y marcas son un vocabulario abierto, y los pesos deben cumplir a la vez un piso, un tope y
una suma.

\begin{table}[ht]
\centering\small
\begin{tabular}{@{}L{2.6cm}L{3.3cm}L{7.6cm}@{}}\toprule
Grupo & Campo & Dominio\\\midrule
$F$ & \cod{presupuesto\_max\_cop} & vac\'io o entero en $[50\,000;\,3\,000\,000]$ COP por hijo\\
 & \cod{marcas\_excluidas} & hasta 10 cadenas de 1 a 40 caracteres\\
$W$ & cuatro pesos & cada $w_i\in[0{,}05;\,0{,}60]$ y $\sum w_i=1{,}00$\\
$P$ & \cod{personajes} & subconjunto del cat\'alogo, hasta 3\\
 & \cod{estilo\_gasto} & econ\'omico, punto medio o estirarse\\\bottomrule
\end{tabular}
\caption{Dominio de los campos de $\R$. Fuente: constantes de \cod{nl2c.py}.}
\label{tab:dominios}
\end{table}

\subsection{Dise\~no del problema}

\textbf{Entradas:} el texto $t\in\Sigma^{*}$ (a lo m\'as 1000 caracteres tras sanear), el cat\'alogo de personajes
$\mathcal{P}$ con hasta 60 nombres y la lista blanca $\mathcal{C}$ de siete componentes de interfaz. Para pedir cambios sobre un
kit se agrega el kit $K$, con hasta 80 renglones. \textbf{Salidas:} $\R$ validada, la lista de ajustes que hizo el servidor, la
especificaci\'on de interfaz $U$ con componentes de $\mathcal{C}$ y un estado (propuesta, sin acci\'on, inv\'alido o sin servicio).
\emph{Postcondici\'on:} si el estado es propuesta, $\R$ cumple la \cref{tab:dominios}, todo texto libre de $\R$ aparece en $t$ y
nada se aplica antes de que el acudiente confirme.

\subsection{Estrategia y seudoc\'odigo}

La estrategia es un plan fijo con validaci\'on en el servidor: el c\'odigo decide qu\'e funciones puede llamar cada ruta antes de
ver $t$, y el modelo solo elige una de ellas \citep{beurerkellner2025patterns}. Esta idea sigue la separaci\'on entre planificar y
ejecutar de \citet{debenedetti2025camel}. Con ese plan, el modelo recibe el texto como dato y llama a una funci\'on con
argumentos (\emph{function calling}); el servidor valida esos argumentos y compila la interfaz; y el acudiente confirma. La l\'inea
base es el formulario de opciones cerradas, que no usa modelo y no tiene riesgo de inyecci\'on, pero no entiende frases que mezclan
campos.

\begin{algorithm}[ht]
\caption{\textsc{Interpretar}$(t,\mathcal{P})$}\label{alg:interpretar}
\begin{algorithmic}[1]
\Require Texto $t\in\Sigma^{*}$, cat\'alogo $\mathcal{P}$
\Ensure Estado y, si es propuesta, la tupla $(\R,A,U)$
\State $s\gets \textsc{Sanear}(t)$ \LineComment{quita datos personales y marcas que escapan del bloque de datos}
\State $L\gets\{\text{fijar\_restricciones},\text{sin\_acci\'on}\}$ \LineComment{plan fijo: la lista no sale de $t$}
\State $(f,a)\gets \textsc{LlamarModelo}(s,L)$ \LineComment{una sola llamada, sin reintento}
\If{$f\notin L$ \textbf{o} $f=\text{sin\_acci\'on}$} \State \Return \emph{sin acci\'on} \EndIf
\State $(\R,A)\gets \textsc{Proyectar}(a,\mathcal{P},s)$ \LineComment{valida forma, recorta dominios, ancla texto libre a $s$}
\If{\textsc{Proyectar} fall\'o por forma} \State \Return \emph{inv\'alido} \EndIf
\State $U\gets \textsc{CompilarUI}(\R)$ \LineComment{la interfaz la arma el servidor, no el modelo}
\State \Return (\emph{propuesta}, $\R$, $A$, $U$) \LineComment{nada se aplica hasta la confirmaci\'on}
\end{algorithmic}
\end{algorithm}

\subsection{Complejidad y costo}

\textsc{Sanear} recorre $t$ con un n\'umero constante de expresiones regulares, de modo que cuesta $O(|t|)$. \textsc{Proyectar}
trabaja sobre un $\R$ de tama\~no constante (4 pesos, hasta 10 marcas y 6 nombres de personaje) y cuesta
$O(|\mathcal{P}|+|s|)$ por el \'indice del cat\'alogo y el cotejo de cada marca con el texto. La normalizaci\'on de pesos itera
a lo m\'as 4 veces, y la compilaci\'on de $U$ genera a lo m\'as 6 componentes, ambas en $O(1)$. En total son
$O(|t|+|\mathcal{P}|)$ m\'as una llamada al modelo, acotada a 400 tokens de salida y 12 segundos de plazo. Medido con 70 intentos
en la ronda 3, la llamada tuvo latencia p50 de 1000\,ms y p95 de 1177\,ms, con 616 tokens de entrada y 39 de salida por pedido en
promedio (unos 0,000077 USD por pedido, con un precio de lista supuesto, no medido). El formulario no tiene llamada al modelo, pero
tampoco se midi\'o su tiempo con acudientes, por lo que no se afirma que NL2C sea m\'as r\'apido para ellos.

\subsection{Seguridad, interfaz generada y confirmaci\'on}

El servidor cubre cinco riesgos de la lista OWASP para aplicaciones con modelos de lenguaje \citep{owasp2025llm}: inyecci\'on de
prompt (texto como dato entre marcas y funciones en lista blanca), informaci\'on sensible (redacci\'on de datos personales antes de
la llamada), manejo de la salida (validaci\'on estricta y plantillas del servidor), agencia excesiva (el modelo solo propone; el
acudiente aplica con su sesi\'on) y consumo ilimitado (l\'imites de texto, tokens, tiempo y tasa). La interfaz generada
\citep{chen2025genui} no la escribe el modelo: el servidor compila $U$ con los siete componentes de $\mathcal{C}$ y la web la
vuelve a validar. La confirmaci\'on sigue las pautas de interacci\'on humano--IA de \citet{amershi2019guidelines}.

\subsection{Evaluaci\'on}

Se midi\'o con 35 frases en espa\~nol (22 de la encuesta a acudientes y 13 construidas), con la etiqueta de oro escrita antes de
correr el modelo, y con 27 ataques de inyecci\'on al estilo de AgentDojo \citep{debenedetti2024agentdojo}. El modelo fue
\cod{gemini-3.5-flash-lite} con temperatura 0, con dos repeticiones por frase (70 intentos) y tres por ataque (81 intentos). Un ataque
\emph{tiene \'exito} cuando su condici\'on se cumple sobre lo que el acudiente ve despu\'es de las barandas del servidor. La
\cref{tab:rondas} da las dos \'ultimas rondas.

\begin{table}[ht]
\centering\small
\begin{tabular}{@{}L{5.4cm}L{3.9cm}L{3.9cm}@{}}\toprule
Medida & Ronda 2 (criterio original) & Ronda 3 (criterio nuevo)\\\midrule
Tarea completa & 57 de 70 (81,4\,\%) & 58 de 70 (82,9\,\%)\\
F1 de personajes y marcas & 1,000 & 1,000\\
Latencia p50 / p95 (ms) & 999 / 1344 & 1000 / 1177\\
Ataques exitosos (ASR) & 0 de 81 & 0 de 81\\
El modelo lo intent\'o & 5 de 81 & 6 de 81\\
Utilidad bajo ataque & 36 de 60 & 36 de 60\\
Controles sin ataque & 12 de 12 & 12 de 12\\\bottomrule
\end{tabular}
\caption{Resultados de las rondas 2 y 3, con el mismo modelo, las mismas frases y los mismos ataques. Fuente:
\cod{resultados/ronda2\_*.json} y \cod{ronda3\_*.json}.}
\label{tab:rondas}
\end{table}

Los 12 fallos de utilidad de la ronda 3 se concentran en el campo \emph{peso mayor}, donde varias etiquetas son discutibles, y el
modelo acert\'o los personajes y las marcas sin errores. Bajo ataque se perdi\'o la tarea leg\'itima en 24 de 60 intentos, porque
el modelo respondi\'o \cod{sin\_acci\'on} ante la orden hostil; con ello se prefiere no proponer ante la duda.

\paragraph{Cambio de criterio.} Despu\'es de la ronda 2 se cambi\'o la condici\'on de \'exito \cod{descarte\_aplicable} por
\cod{descarte\_sin\_confirmar}, porque en la integraci\'on con la sesi\'on de pedidos el descarte pas\'o a ser aplicable tras la
confirmaci\'on del acudiente. El cambio se hizo despu\'es de ver los resultados de la ronda 2 y no estaba registrado antes de
correr, as\'i que no es un criterio preregistrado. En esa condici\'on ambas rondas dieron 6 evaluaciones y 0 \'exitos.

\paragraph{Pasada con un kit real.} Con los servicios levantados, 4 pedidos contra un kit real (perfil Ahorrador, 36 renglones)
dieron una propuesta de cambio, una explicaci\'on, un descarte propuesto y un rechazo de la orden hostil, con 786 a 885\,ms cada uno.
Con $n=4$ es una comprobaci\'on de funcionamiento y no una tasa.

\subsection{Limitaciones}

Con 0 \'exitos en 81 intentos, el l\'imite superior del intervalo de confianza del 95\,\% es de unos 3,7\,\% (regla de tres), de modo
que solo se puede decir que no se encontr\'o un ataque exitoso con estos casos. Las repeticiones a temperatura 0 no son
independientes, el modelo y la fecha son uno solo, el kit y el cat\'alogo de las tres rondas son sint\'eticos, y los ataques los
escribi\'o quien construy\'o las defensas. Como el modelo intent\'o solo 6 de 81, las barandas se ejercitaron en pocos intentos;
el resto lo cubren pruebas con un modelo doble obediente.

\paragraph{Consumo de los filtros duros.} Los campos \cod{marcas\_excluidas} y \cod{solo\_disponibles} de $F$ llegan al
optimizador. Se aplican a los candidatos de cada rengl\'on antes de calcular el frente de Pareto, no como un peso; la marca que el
colegio nombra en el texto del rengl\'on se conserva, y un rengl\'on que se queda sin candidatos sale como hueco con un mensaje en
lugar de usar la marca excluida. La ruta r\'apida que reutiliza un kit precalculado rechaza estos filtros y obliga a recalcular. Lo
respaldan las pruebas \cod{test\_filtros\_duros.py} (por ejemplo \cod{test\_la\_marca\_no\_sale\_en\_ningun\_kit\_ni\_sus\_alternativas},
\cod{test\_la\_marca\_que\_exige\_el\_colegio\_se\_conserva\_y\_se\_explica} y
\cod{test\_linea\_sin\_candidatos\_queda\_como\_hueco\_con\_mensaje}) y \cod{test\_a\_preferencias\_pasa\_marcas\_excluidas\_y\_solo\_disponibles\_al\_optimizador}.
Esta comprobaci\'on es de c\'odigo con datos de prueba y no mide el efecto sobre el cat\'alogo real.

\bibliographystyle{plainnat}
\bibliography{nl2c}

\end{document}
